\section{Conformit\`a alla consegna}
\label{sec:compliance}

La Tab.~\ref{tab:compliance} mappa ogni requisito della consegna sul relativo
deliverable. Legenda stato: \checkmark{} soddisfatto, \textasciitilde{}
parziale.

\begin{table}[H]
\centering
\small
\caption{Matrice di conformit\`a requisiti $\to$ implementazione.}
\label{tab:compliance}
\begin{tabular}{p{6.6cm}cp{5.2cm}}
\toprule
Requisito & Stato & Evidenza \\
\midrule
\multicolumn{3}{l}{\textit{Env 1 --- CartPole (comune)}} \\
Q\hbox{-}Learning, SARSA, MC, TD(0) ``from scratch'' & \checkmark &
  \code{rl\_lab/algorithms/tabular/} \\
Discretizzazione stato continuo & \checkmark &
  \code{envs/discretizer.py}; Tab.~\ref{tab:envs} \\
Curve di apprendimento (reward vs episodi) & \checkmark &
  Fig.~\ref{fig:cp-control} \\
Sensibilit\`a a $\varepsilon,\alpha,\gamma$ & \checkmark &
  Sez.~\ref{sec:res-cartpole-sens}, 6 figure \\
SARSA vs Q\hbox{-}Learning (empirico + teorico) & \checkmark &
  Sez.~\ref{sec:res-cartpole-control}, \ref{sec:theory} \\
MC vs TD(0) (bias/varianza/stabilit\`a) & \checkmark &
  Fig.~\ref{fig:pred-value} (stima $V$); vedi fix \#1 in Sez.~\ref{sec:lessons-fixes} \\
$\geq 1$ caso di fallimento & \checkmark &
  Sez.~\ref{sec:res-cartpole-fail}, 2 casi \\
\midrule
\multicolumn{3}{l}{\textit{Env 2 --- Acrobot (assegnato)}} \\
Stessi algoritmi applicati & \checkmark &
  \code{run\_acrobot.py} \\
Difficolt\`a di esplorazione (reward sparso) & \checkmark &
  Sez.~\ref{sec:res-acrobot-eps} \\
Necessit\`a di discretizzazione/approssimazione & \checkmark &
  Sez.~\ref{sec:res-acrobot-bins} (esperimento sui bin) \\
\midrule
\multicolumn{3}{l}{\textit{Env 3 --- Pong (Atari, assegnato)}} \\
DQN ``from scratch'' (replay + target) & \checkmark &
  \code{algorithms/deep/dqn.py} \\
Ablazione con/senza target network & \checkmark &
  Fig.~\ref{fig:pong}, Tab.~\ref{tab:pong-ablation} \\
Ablazione con/senza experience replay & \checkmark &
  Fig.~\ref{fig:pong}, Tab.~\ref{tab:pong-ablation} \\
Apprendimento vs policy casuale & \checkmark &
  baseline $\approx-20.4$; DQN $-12.1$ \\
Double DQN (bonus) & \checkmark &
  flag \code{use\_double\_dqn}; Sez.~\ref{sec:res-pong} \\
\midrule
\multicolumn{3}{l}{\textit{Requisiti trasversali}} \\
Media su 3 seed + deviazione standard & \checkmark &
  tutte le figure/tabelle \\
Riproducibilit\`a & \checkmark &
  seeding, \code{.npz}, \code{uv.lock} \\
No librerie RL di alto livello (Env 1/2) & \checkmark &
  core scritto a mano; SB3 non usato \\
Report 10--15 pagine, con failure analysis & \checkmark &
  questo documento (LaTeX) \\
\bottomrule
\end{tabular}
\end{table}

\paragraph{Nota sul confronto MC vs TD(0).}
In una prima versione questo era l'unico requisito solo parzialmente
soddisfatto: la visualizzazione mostrava il return della policy (identico per
costruzione) e non il compromesso bias/varianza. Con il fix \#1
(Sez.~\ref{sec:lessons-fixes}) si traccia ora anche la stima $V(s_{\text{ref}})$
episodio per episodio, che espone la varianza alta di MC contro la traiettoria
liscia di TD(0): il requisito \`e ora pienamente soddisfatto. Resta come
possibile estensione una misura globale (RMSE di $V$ su tutti gli stati).
